
\section{Introduction}
\label{sec:introduction}

% Distributed systems rely on multiple interconnected nodes that communicate and cooperate to provide services.
% In this kind of environment, node availability is important because a node may crash, become unreachable,
% respond slowly, or lose connection with other nodes.
% When these problems are not detected and handled properly, they can affect communication, reduce
% reliability, and weaken the overall operation of the system.

% Node health checking is used to observe whether nodes remain active and responsive.
% In many systems, this responsibility is assigned to a central monitoring component because it is easier to
% collect, manage, and visualize node status from one place.
% However, this structure can introduce a single point of failure when the central component becomes
% unavailable. This project therefore focuses on a decentralized node health-checking structure, where
% monitoring responsibility is distributed among the nodes themselves instead of being concentrated in one component.

% In the proposed system, each node maintains relationships with a limited number of neighboring nodes.
% These relationships define which nodes are responsible for monitoring one another and form the basic
% structure of the decentralized health-checking mechanism.
% Through node-to-node communication, the system identifies healthy, suspicious, and failed nodes.
% When a node or connection problem is detected, the system updates the affected relationships and attempts
% to preserve the monitoring structure so that health checking can continue.

Node health checking is a fundamental requirement in distributed systems because the timely detection of
failed or unreachable nodes directly affects service continuity, coordination correctness, and system
reconfiguration \cite{liEnergyefficientLazyGroup2025} \cite{yangMedleyMembershipService2022}.
The impact of an undetected node failure can be severe, particularly when the failed node is topologically or
functionally critical; in interdependent networks, even a small fraction of node failures can trigger
recursive cascades and widespread fragmentation \cite{buldyrevCatastrophicCascadeFailures2010}.
A concrete real-world example is the 2003 FirstEnergy alarm system failure, in which a critical software node
stalled silently due to an unmonitored race condition, blinding operators and allowing a localized grid
anomaly to propagate into a catastrophic, multi-regional cascading blackout \cite{Blackout20032003}.

A common health-checking paradigm relies on centralized supervision, in which a server or a small set of
servers monitors nodes through periodic liveness messages or heartbeat-style reporting
\cite{vanrenesseGossipStyleFailureDetection1998} \cite{kawazoeaguileraHeartbeatTimeoutfreeFailure1997}.
Heartbeat-based detection remains widely used because timeout on missing periodic messages provides a simple
operational rule for identifying suspected failures \cite{vanrenesseGossipStyleFailureDetection1998}
\cite{kawazoeaguileraHeartbeatTimeoutfreeFailure1997}.
Such designs are also attractive from an engineering perspective because centralized control simplifies
operational management and supports a more coherent view of membership and monitoring state;
distributed-systems research continues to build many higher-level services on top of membership knowledge
\cite{huntZooKeeperWaitfreeCoordination}.

However, centralized monitoring does not scale uniformly well as the system grows.
Large-scale and geographically distributed deployments make manual investigation difficult, which is why
automatic failure detection becomes essential in IoT and other distributed environments
\cite{yangMedleyMembershipService2022} \cite{banerjeeHardwareAssistedHeartbeatMechanism2022}.
In addition, centralized streaming and monitoring architectures can encounter bandwidth and concurrency
bottlenecks, along with reduced fault tolerance under heavy load
\cite{guoDecentralizedMultiVenueRealTime2025}
\cite{weiSwarmCostEfficientVideo2024}.
These limitations motivate decentralized alternatives.

In decentralized health checking, nodes monitor one another without relying on a single global coordinator
\cite{yangMedleyMembershipService2022}
\cite{nikolicSelfHealingDilemmasDistributed2021}.
% Keywords: decentralized, single, health, universal, central
% Nikolic
% 2020 --> feel indirectly implying
% Yang
% 2022
This design improves robustness and reduces dependence on central infrastructure, but it also introduces
harder problems in membership maintenance, coordination, and incomplete local knowledge
\cite{nikolicSelfHealingDilemmasDistributed2021}.
% Keywords: Decentralized, byzantine, malicious, incorrect, large-scale
% Nikolic
% 2020
% Rodr'iguez
% 2020 --> can't see citations/ referneces
Within this line of work, three recurring research concerns are network formation or neighbor discovery,
failure detection accuracy and efficiency, and fault tolerance or self-healing after a node failure
\cite{rodriguezDecentralisedSelfHealingApproach2021}.
% Rodr'iguez
% 2020

A representative decentralized protocol is SWIM, which will be discussed at length in
section~\ref{sec:related_work} \cite{dasSWIMScalableWeaklyconsistent2002}.
Its main advantage over traditional heartbeating is scalability: SWIM separates failure detection from
membership dissemination, and both the expected detection time and expected message load per member remain
independent of group size \cite{dasSWIMScalableWeaklyconsistent2002}.
Related systems have extended this direction to specialized environments. For example, Medley adapts
decentralized membership and failure detection to wireless IoT settings through spatially biased probing,
reducing the product of detection time and communication cost while maintaining low false positive rates
\cite{yangMedleyMembershipService2022}.
Other work has shown that decentralized self-healing can rely on locally maintained neighborhood topology
information to recreate and reconnect failed nodes after detection \cite{rodriguezDecentralisedSelfHealingApproach2021}.

% \subsection{Background}
% \label{subsec:background}

% A node health-checking system observes the condition of nodes in a distributed environment.
% A node may be considered healthy when it responds to health-check messages within an acceptable time, and
% unhealthy when it becomes unreachable, stops responding, or behaves abnormally due to system or network issues.
% Health checking is therefore related to failure detection, connection monitoring, and membership management.

% In a decentralized health-checking structure, nodes do not need a complete view of the entire system.
% Instead, each node is responsible for monitoring a smaller set of neighboring nodes.
% This reduces the amount of direct monitoring required from each participant and allows the health-checking
% responsibility to be shared across the network.
% The structure of node relationships is therefore important because it determines how monitoring tasks are
% distributed and how the system reacts when nodes join, leave, fail, or recover.

% This project uses a neighbor-based monitoring structure.
% Each node attempts to maintain a fixed number of neighbors, denoted as \textbf{k}.
% The value of \textbf{k} represents the desired number of monitoring relationships for each node. When this
% number is reduced because of failure or disconnection, the system attempts to repair the missing
% relationships by connecting to other available nodes.
% In this way, the health-checking structure is not static, but can adapt to changes in the network.

% \subsection{Motivation}
% \label{subsec:motivation}

% The motivation of this project is to design a health-checking mechanism that remains active through local
% node cooperation.
% Instead of treating monitoring as a separate external service, the system makes health checking part of the
% behavior of each node.
% This allows nodes to participate directly in detecting failures, maintaining relationships, and restoring
% the monitoring structure when changes occur.

% A decentralized structure is useful when the system must continue operating under partial failure conditions.
% In such conditions, some nodes or connections may fail while others remain functional.
% The health-checking mechanism should therefore avoid depending on a complete and perfect network view. It
% should instead allow each node to make local decisions based on the neighbors it can observe.

% This project is motivated by the need to combine failure detection with relationship maintenance.
% Detecting that a node is unavailable is not enough by itself; the system must also respond by updating the
% monitoring topology.
% For this reason, the design focuses not only on identifying unhealthy nodes, but also on preserving the
% network structure through joining, repair, and rewiring mechanisms.

% \subsection{Contributions}
% \label{subsec:contribution}

% This project presents a decentralized node health-checking system in which nodes monitor one another
% through neighbor-based relationships.
% The main contributions are as follows:

% \begin{itemize}

% \item \textbf{Neighbor-based decentralized structure:}
% The system distributes health-checking responsibility across nodes by assigning each node a limited set of
% neighbors to monitor.

% \item \textbf{Fixed-\(k\) relationship maintenance:}
% Each node attempts to maintain a fixed number of neighbor relationships, allowing the monitoring structure
% to remain organized as the network changes.

% \item \textbf{Node joining mechanism:}
% The system supports new nodes entering the network and becoming part of the decentralized monitoring topology.

% \item \textbf{Failure detection mechanism:}
% Nodes use health-check communication to identify healthy, suspicious, and failed neighbors based on their
% observed responses.

% \item \textbf{Repair and rewiring mechanism:}
% When nodes or connections become unavailable, the system attempts to replace missing relationships and
% restore the required monitoring structure.

% \item \textbf{Optional observation layer:}
% A dashboard, server, and database are included to display node status and failure events, but they are
% separate from the core decentralized health-checking mechanism.

% \item \textbf{Experimental evaluation:}
% The system is tested under normal operation, node joining, node crash, temporary disconnection, unstable
% connection, recovery, and observation-layer unavailability scenarios.

% \end{itemize}

% Overall, this project contributes a working prototype of a decentralized node health-checking system that
% combines local monitoring, relationship maintenance, failure detection, and topology repair.
% The system demonstrates how node cooperation can be used to maintain health-checking functionality in a
% changing distributed environment.
